\documentclass[journal]{IEEEtran}

\usepackage{amsmath,amssymb}
\usepackage{graphicx}
\usepackage{booktabs}
\usepackage{multirow}
\usepackage{cite}
\usepackage{microtype}
\usepackage{url}
\usepackage[hidelinks]{hyperref}

\interdisplaylinepenalty=2500

\title{Green-PLUM: Energy-Aware Bidirectional Bandwidth Coordination for Video Conferencing}

\author{Author Name(s)}

\begin{document}

\maketitle

\begin{abstract}
This document is the initial manuscript skeleton for a planned submission to
\emph{IEEE/ACM Transactions on Networking}. It extends PLUM with energy- and
thermal-aware bidirectional bandwidth coordination for mobile video
conferencing. Replace this placeholder abstract after the problem statement,
method, and evaluation scope are finalized.
\end{abstract}

\begin{IEEEkeywords}
Video conferencing, bidirectional bandwidth coordination, energy efficiency,
thermal management, quality of experience.
\end{IEEEkeywords}

\section{Introduction}
\IEEEPARstart{M}{obile} video conferencing must jointly manage uplink and
downlink traffic under shared or half-duplex bottlenecks. PLUM established the
baseline coordination framework~\cite{plum2024}. This manuscript studies how
device battery state and thermal constraints can be incorporated without
discarding PLUM's network-level benefits.

\section{Motivation and Problem Statement}
% TODO: Explain the limitations of QoE-only allocation and define the target
% operating scenarios, device state, and optimization objectives.

\section{System Model}
% TODO: Describe the SFU architecture, bottleneck model, rate variables,
% battery model, power model, and thermal dynamics.

\section{Green-PLUM Design}
% TODO: Present the objective, constraints, solver protocol, state updates,
% and compatibility with the original PLUM control loop.

\section{Implementation}
% TODO: Document the NS-3.37 videoconf module changes, Python/SLSQP solver,
% trace-driven setup, and experiment flags.

\section{Evaluation}
% TODO: Separate standalone model results, model-based Wi-Fi experiments, and
% trace-driven packet-level experiments. Do not mix their claims.

\begin{table}[t]
  \caption{Placeholder for the Main Evaluation Summary}
  \label{tab:summary}
  \centering
  \begin{tabular}{lccc}
    \toprule
    Policy & Network QoE & Energy & Survival \\
    \midrule
    Vanilla & TBD & TBD & TBD \\
    PLUM & TBD & TBD & TBD \\
    Green-PLUM & TBD & TBD & TBD \\
    \bottomrule
  \end{tabular}
\end{table}

\section{Related Work}
% TODO: Cover bidirectional coordination, mobile video energy measurement,
% codec/network power models, and thermal-aware systems.

\section{Discussion and Limitations}
% TODO: State which energy values are literature-calibrated models rather than
% measurements, and discuss the QoE--energy tradeoff and generalizability.

\section{Conclusion}
% TODO: Summarize only claims supported by the final validated experiments.

\bibliographystyle{IEEEtran}
\bibliography{references}

\end{document}
